\documentclass[11pt,reqno]{amsart}

\usepackage[T1]{fontenc}
\usepackage{amsmath,amssymb,amsthm}
\usepackage[margin=1.15in]{geometry}
\usepackage{booktabs}
\usepackage[colorlinks=true,linkcolor=black,citecolor=black,urlcolor=blue]{hyperref}

\theoremstyle{plain}
\newtheorem{theorem}{Theorem}[section]
\newtheorem{proposition}[theorem]{Proposition}
\newtheorem{lemma}[theorem]{Lemma}
\newtheorem{corollary}[theorem]{Corollary}
\theoremstyle{definition}
\newtheorem{definition}[theorem]{Definition}
\newtheorem{remark}[theorem]{Remark}
\newtheorem{observation}[theorem]{Observation}

\newcommand{\Z}{\mathbb{Z}}
\newcommand{\N}{\mathbb{N}}
\newcommand{\Ll}{\mathsf{L}}
\newcommand{\Rr}{\mathsf{R}}
\newcommand{\Mon}{\mathcal{M}}
\newcommand{\Nn}{\mathcal{N}}
\newcommand{\tr}{\operatorname{tr}}
\newcommand{\SL}{\mathrm{SL}}
\newcommand{\GL}{\mathrm{GL}}
\newcommand{\sqm}[4]{\left(\begin{smallmatrix}#1&#2\\#3&#4\end{smallmatrix}\right)}

\begin{document}

\title[Necklaces and Zagier-reduced forms]
      {Necklaces, the letter swap, and Zagier-reduced forms of discriminant $N^2-4$}

\author{}
\date{12 August 2026}

\begin{abstract}
Let $\Mon$ be the free monoid generated by
$\Ll=\sqm{1}{1}{0}{1}$ and $\Rr=\sqm{1}{0}{1}{1}$, and let
$\Phi(N)=\#\{X\in\Mon:\tr X=N\}$ be the state count of the Farey fraction spin chain of
Kleban and \"Ozl\"uk. Let $Z(D)$ denote the number of Zagier-reduced integral binary
quadratic forms of discriminant~$D$. The Online Encyclopedia of Integer Sequences has
recorded since 2015 the conjecture that $\Phi(N)=2Z(N^2-4)$
(\texttt{A264598}, \texttt{A264597}, \texttt{A257007}), with the annotation
``It would be nice to have a proof''. We prove it, in three ways.

First, an explicit divisor parametrisation of the Zagier-reduced forms of a non-square
discriminant $D$ (Lemma~\ref{lem:param}) identifies $Z(D)$ with the quantity
$\Upsilon^{<}(D)$ appearing as an intermediate step in Technau's 2023 analysis of
$\Phi$; the conjecture is then immediate from his equation~(3.2). Second, the same
parametrisation combined with a fixed-point-free involution on divisor pairs yields a
short self-contained proof from first principles, independent of the literature. Third,
and structurally, we bend the words of $\Mon$ into necklaces: cyclic words of trace $N$
are the conjugacy classes of trace $N$, the length of the Zagier cycle of a class equals
the number of $\Rr$'s in its word, and conjugation by $\sqm{0}{1}{1}{0}$ --- a matrix of
determinant $-1$, hence not a symmetry of $\SL_2(\Z)$, but a trace-preserving one ---
permutes the classes of trace $N$ while exchanging the two letter counts. The factor $2$
in the conjecture is exactly the two colours of bead. This third argument proves the
finer, per-class statement of which the conjecture is the total.

We then introduce the \emph{cut-sum} $S(w)$ of a necklace, the sum of the matrices of all
its cuts: a well-defined matrix-valued invariant strictly finer than the trace. We prove
$\det S=n+\sum \tr[u,v]$, summed over pairs of cuts with $u,v$ the two arcs they
determine, and identify $(\tr S)^2-4\det S$ as the Minkowski norm of a sum of pole
vectors. Finally we record the total-matrix identity
$\sum_{\tr X=N}X=\tfrac{N}{2}\Phi(N)\,I+\bigl(\sum_{a}\sigma(a(N-a)-1)\bigr)J$, whose
$\sigma$-half appears not to have been studied. All assertions were verified by exact
integer computation.
\end{abstract}

\maketitle

\section{Introduction}

\subsection{The spin chain}
Let
\[
  \Ll=\begin{pmatrix}1&1\\0&1\end{pmatrix},\qquad
  \Rr=\begin{pmatrix}1&0\\1&1\end{pmatrix},
\]
and let $\Mon\subset\SL_2(\Z)$ be the monoid they generate. Kleban and \"Ozl\"uk
\cite{KO} introduced a statistical-mechanical model --- the \emph{Farey fraction spin
chain} --- whose states are the elements of $\Mon$ and whose energy is the logarithm of
the trace. The associated counting function is
\begin{equation}\label{eq:Phi}
  \Phi(N)\;=\;\#\{X\in\Mon:\tr X=N\},\qquad N\ge 3 .
\end{equation}
It begins $2,6,8,14,14,24,16,40,26,36,\dots$ and is \texttt{A264598} in the OEIS
\cite{OEISA264598}. Its behaviour is notoriously irregular: Kleban and \"Ozl\"uk
conjectured $\Phi(N)\sim c_0N\log N$, and Peter \cite{Peter} showed that no such
asymptotic can hold, since $\Phi(N)/(N\log N)$ possesses a smooth limiting distribution
rather than a limit. Attention has since concentrated on the summatory function, treated
by Kallies--\"Ozl\"uk--Peter--Snyder \cite{KOPS}, Boca \cite{Boca} and Ustinov
\cite{Ustinov}. Most recently Technau \cite{Technau} reduced the estimation of $\Phi$ to
a divisor problem for quadratic polynomials solved by Hooley \cite{Hooley} and, in full
generality, by Bykovski\u\i--Ustinov \cite{BU}.

\subsection{Zagier-reduced forms}
Independently, Zagier's reduction theory for indefinite binary quadratic forms
\cite[\S13]{Zagier} calls an integral form $Q=(A,B,C)$, meaning $Ax^2+Bxy+Cy^2$, of
discriminant $D=B^2-4AC>0$ \emph{reduced} when
\begin{equation}\label{eq:zagred}
  A\ge 1,\qquad C\ge 1,\qquad B>A+C .
\end{equation}
Write $Z(D)$ for the number of such forms. The values of $Z(N^2-4)$ for $N=3,4,5,\dots$
are $1,3,4,7,7,12,8,20,13,18,\dots$, the sequence \texttt{A257007} \cite{OEISA257007}.

\subsection{The conjecture}
Comparing the two lists, Sloane observed in November 2015 that each term of
\texttt{A264598} is twice the corresponding term of \texttt{A257007}, and recorded in
\texttt{A264598} the remark ``\emph{This also appears to be the same as \texttt{A257007}
multiplied by $2$. It would be nice to have a proof.}'' The same caveat is repeated in
\texttt{A257007}, whose \textsc{Maple} program carries the warning that it should not be
used to extend the sequence until the identity is proved. Both annotations stand today.
The present note removes them.

\begin{theorem}\label{thm:main}
For every integer $N\ge 3$,
\[
  \Phi(N)\;=\;2\,Z(N^2-4).
\]
\end{theorem}

\subsection{Three proofs, and what is new in each}
We give three arguments, and we are explicit about how much of each is already in the
literature.

\emph{(i) A deduction from \cite{Technau}.} Technau's route to his Theorem~2 passes
through the intermediate identity $\Phi(N)=2\Upsilon^{<}(N^2-4)$ \cite[(3.2)]{Technau},
where $\Upsilon^{<}$ is a certain restricted divisor sum. Our Lemma~\ref{lem:param}
below shows that $\Upsilon^{<}(D)$ is precisely the number of Zagier-reduced forms of
discriminant $D$; Theorem~\ref{thm:main} follows at once. The factor $2$ was therefore
already proved in 2023. What was missing was the dictionary entry: \cite{Technau} does
not mention Zagier reduction, and the OEIS entries do not mention \cite{Technau}. The
new content of this proof is Lemma~\ref{lem:param}.

\emph{(ii) An elementary self-contained proof} (\S\ref{sec:involution}). Combining
Lemma~\ref{lem:param} with a fixed-point-free involution on divisor pairs gives a proof
of Theorem~\ref{thm:main} occupying half a page and using nothing beyond the freeness of
$\Mon$. The involution is a re-derivation, in the coordinates natural for Zagier
reduction, of \cite[Lemma~4]{Technau}.

\emph{(iii) A structural proof} (\S\ref{sec:necklace}), which is the reason for this
note. Bend the words of $\Mon$ into loops. Cyclic words --- \emph{necklaces} --- of
trace $N$ are exactly the $\SL_2(\Z)$-conjugacy classes of trace $N$; each contributes
to $\Phi(N)$ once for each of its cuts, so
$\Phi(N)=\sum_{\mathcal{C}}(\#\Ll+\#\Rr)$. On the other side, the length of the Zagier
cycle of a class equals $\#\Rr$ alone (Proposition~\ref{prop:zagierR}). The two are
reconciled by the involution induced by conjugation with $J=\sqm{0}{1}{1}{0}$, which
interchanges $\Ll$ and $\Rr$: it has determinant $-1$, so it is not an inner symmetry of
$\SL_2(\Z)$, but it preserves trace and determinant, hence permutes the classes of trace
$N$ among themselves while exchanging the two letter counts. Summed over a set that it
permutes, the counts must agree. This proof establishes the finer per-class statement,
of which Theorem~\ref{thm:main} is the total.

\subsection{Further results}
Section~\ref{sec:swapinv} determines the necklaces fixed by the letter swap: they are
exactly the words $u\bar u$, and their traces are exactly the integers $s^2+2$.
Section~\ref{sec:cutsum} introduces the cut-sum, a matrix-valued necklace invariant
finer than the trace, and proves a commutator formula for its determinant together with
a Minkowski interpretation of its discriminant. Section~\ref{sec:total} records a
total-matrix identity whose off-diagonal entry is a sum of $\sigma(a(N-a)-1)$, a
quantity that seems not to have been considered. Section~\ref{sec:verify} documents the
computations, and Section~\ref{sec:open} lists what we could not settle.

\medskip
\noindent\emph{Conventions.} $d(\cdot)$ and $\sigma(\cdot)$ are the number-of-divisors
and sum-of-divisors functions. Divisors are positive. Discriminants are positive and
never perfect squares; note $N^2-4$ is a square only for $N=2$.

\section{The free monoid and its trace fibres}\label{sec:monoid}

\begin{proposition}\label{prop:free}
$\Mon\cup\{I\}$ is precisely the set of $2\times2$ integer matrices with non-negative
entries and determinant $1$, and it is free on $\{\Ll,\Rr\}$.
\end{proposition}

\begin{proof}
Freeness is due to Kleban and \"Ozl\"uk \cite{KO}; see also Nathanson \cite{Nathanson}
for a proof of a more general statement by the ping-pong lemma.

For surjectivity, let $X=\sqm{a}{b}{c}{d}$ have non-negative entries and $ad-bc=1$, and
suppose $X\neq I$. We claim that either $a\ge c$ and $b\ge d$, or $c\ge a$ and $d\ge b$.
Suppose not; then one of the following holds.
\begin{itemize}
\item $a<c$ and $d<b$. All entries being non-negative integers, $c\ge a+1$ and
$b\ge d+1$, so $bc\ge(a+1)(d+1)=ad+a+d+1>ad$, contradicting $ad-bc=1$.
\item $c<a$ and $b<d$. Symmetrically $ad\ge(b+1)(c+1)=bc+b+c+1$, so
$1=ad-bc\ge b+c+1$, forcing $b=c=0$ and then $ad=1$, i.e.\ $X=I$, which was excluded.
\end{itemize}
(The two remaining combinations, $a<c$ with $c<a$ and $b<d$ with $d<b$, are vacuous.)
In the first case $\Ll^{-1}X=\sqm{a-c}{b-d}{c}{d}$ again has non-negative entries and
determinant~$1$; in the second, $\Rr^{-1}X=\sqm{a}{b}{c-a}{d-b}$ does. Since the sum of
the entries strictly decreases at each such step unless $X=I$, iteration expresses $X$
as a word in $\Ll$ and $\Rr$.
\end{proof}

\begin{proposition}\label{prop:fibre}
Let $N\ge3$ and put $\Nn_a=a(N-a)-1$. The map $X\mapsto(X_{11},X_{12})$ is a bijection
\[
  \{X\in\Mon:\tr X=N\}\;\longrightarrow\;
  \bigl\{(a,b)\in\Z^2:\ 1\le a\le N-1,\ b\ge1,\ b\mid \Nn_a\bigr\}.
\]
Consequently
\begin{equation}\label{eq:Phidiv}
  \Phi(N)\;=\;\sum_{a=1}^{N-1}d\bigl(a(N-a)-1\bigr).
\end{equation}
\end{proposition}

\begin{proof}
Write $X=\sqm{a}{b}{c}{d}$ with non-negative entries, $a+d=N$, $ad-bc=1$; then
$bc=ad-1=\Nn_a$. If $a=0$ or $a=N$ then $\Nn_a=-1<0$, impossible for non-negative
$b,c$; so $1\le a\le N-1$ and $\Nn_a\ge (N-1)-1\ge1$. Hence $b\ge1$, $b\mid \Nn_a$, and
$X$ is recovered from $(a,b)$ by $d=N-a$, $c=\Nn_a/b$. Conversely every such pair
produces a matrix with non-negative entries, determinant~$1$ and trace~$N$, which lies
in $\Mon$ by Proposition~\ref{prop:free} (it is not $I$ since $N\ge3$).
\end{proof}

\begin{remark}\label{rem:technau2}
Identity \eqref{eq:Phidiv} is Technau's Theorem~2 \cite{Technau} in different
coordinates. Indeed, put $m=N-2a$; as $a$ runs over $1,\dots,N-1$, $m$ runs over exactly
those integers with $m^2<N^2-4$ and $m^2\equiv N^2-4\pmod 4$, and
$\bigl((N^2-4)-m^2\bigr)/4=a(N-a)-1$. Thus \eqref{eq:Phidiv} reads
$\Phi(N)=\Upsilon(N^2-4)$ with
$\Upsilon(X)=\sum_{m^2<X,\;m^2\equiv X (4)}d\bigl((X-m^2)/4\bigr)$, which is
\cite[Theorem~2]{Technau}. We include the short proof above to keep the present note
self-contained.
\end{remark}

\section{A divisor parametrisation of the Zagier-reduced forms}\label{sec:param}

Throughout this section $D>0$ is an integer, not a perfect square, with
$D\equiv0,1\pmod4$. Put
\begin{equation}\label{eq:P}
  P(D)=\Bigl\{(h,k)\in\Z^2:\ k^2<D,\ k^2\equiv D\!\!\pmod 4,\ h\ge1,\
  h\mid\tfrac{D-k^2}{4}\Bigr\},
\end{equation}
and for $(h,k)\in P(D)$ write $g=g(h,k)=\frac{D-k^2}{4h}\ge1$ for the complementary
divisor, so that $4hg=D-k^2$. Let
\begin{equation}\label{eq:Pplus}
  P^{+}(D)=\bigl\{(h,k)\in P(D):\ 2h+k>\sqrt{D}\,\bigr\}.
\end{equation}

\begin{lemma}\label{lem:param}
The assignment $(A,B,C)\mapsto(h,k):=(A,\,B-2A)$ is a bijection from the set of
Zagier-reduced forms of discriminant $D$ onto $P^{+}(D)$, with inverse
$(h,k)\mapsto(A,B,C)=\bigl(h,\;k+2h,\;h+k-g(h,k)\bigr)$. In particular
\[
  Z(D)=\#P^{+}(D).
\]
\end{lemma}

\begin{proof}
Let $(A,B,C)$ satisfy \eqref{eq:zagred} with $B^2-4AC=D$, and set $h=A\ge1$,
$k=B-2A$. Then
\[
  D=B^2-4AC=(k+2h)^2-4hC=k^2+4h\,(h+k-C),
\]
so with $g:=h+k-C$ we have $4hg=D-k^2$; in particular $k^2\equiv D\pmod4$ and, provided
$g\ge1$, also $h\mid (D-k^2)/4$ and $k^2<D$. Now
\[
  B>A+C \iff k+2h>h+C=h+(h+k-g) \iff g\ge 1 ,
\]
and $g\ge1$ is equivalent to $k^2<D$ because $4hg=D-k^2$ and $h\ge1$. Next, since
$h\ge1$,
\[
  (2h+k)^2>D \iff 4h^2+4hk+k^2>4hg+k^2 \iff h+k>g \iff C=h+k-g\ge1 .
\]
If $C\ge1$ then $h+k>g\ge1$, so $2h+k=h+(h+k)>0$ and therefore $2h+k>\sqrt D$; conversely
$2h+k>\sqrt D>0$ gives $(2h+k)^2>D$ and hence $C\ge1$. Thus the constraints
\eqref{eq:zagred} correspond exactly to membership in $P^{+}(D)$. The displayed inverse
map is immediate, and both composites are the identity.
\end{proof}

\begin{corollary}\label{cor:upsilon}
$Z(D)=\Upsilon^{<}(D)$, where $\Upsilon^{<}$ is the restricted divisor sum of
\cite[(3.3)]{Technau}.
\end{corollary}

\begin{proof}
By definition \cite[(3.3)]{Technau},
$\Upsilon^{<}(D)=\sum_{m^2<D,\,m^2\equiv D(4)}d_{(m+\sqrt D)/2}\bigl((D-m^2)/4\bigr)$,
where $d_n(x)=\#\{k\mid x: k<n\}$; that is, $\Upsilon^{<}(D)$ counts the pairs $(h,m)$
with $h\mid (D-m^2)/4$ and $2h-m<\sqrt D$. Substituting $k=-m$, which permutes the index
set, $\Upsilon^{<}(D)=\#\{(h,k)\in P(D): 2h+k<\sqrt D\}=\#P(D)-\#P^{+}(D)$; equality is
excluded because $D$ is not a square. By \cite[Lemma~4]{Technau},
$\Upsilon(D)=\#P(D)=2\Upsilon^{<}(D)$ for non-square $D$, whence
$\#P^{+}(D)=\Upsilon^{<}(D)$, and Lemma~\ref{lem:param} concludes.
\end{proof}

\begin{proof}[First proof of Theorem~\ref{thm:main}]
By \cite[(3.2)]{Technau}, $\Phi(N)=2\Upsilon^{<}(N^2-4)$ for $N\ge3$, and
$N^2-4$ is not a square. Corollary~\ref{cor:upsilon} identifies $\Upsilon^{<}(N^2-4)$
with $Z(N^2-4)$.
\end{proof}

\section{An elementary proof}\label{sec:involution}

The following is independent of \cite{Technau} and of all reduction theory; it uses only
Proposition~\ref{prop:fibre} and Lemma~\ref{lem:param}.

\begin{lemma}\label{lem:invol}
Let $D>0$ be a non-square with $D\equiv0,1\pmod4$. The map
\[
  \iota(h,k)=\bigl(g(h,k),\,-k\bigr),\qquad g(h,k)=\tfrac{D-k^2}{4h},
\]
is a fixed-point-free involution of $P(D)$ which maps $P^{+}(D)$ onto its complement.
Consequently $\#P^{+}(D)=\tfrac12\#P(D)$.
\end{lemma}

\begin{proof}
Since $4hg=D-k^2=D-(-k)^2$ and $g\ge1$ is an integer dividing $(D-k^2)/4$ with
complementary divisor $h$, the map $\iota$ sends $P(D)$ to itself, and
$\iota\circ\iota=\mathrm{id}$.

A fixed point would require $k=-k$ and $h=g$, i.e.\ $k=0$ and $D=4h^2$, contradicting
that $D$ is not a square.

By the computation in the proof of Lemma~\ref{lem:param}, $(h,k)\in P^{+}(D)$ if and
only if $h+k>g$, and likewise $\iota(h,k)=(g,-k)\in P^{+}(D)$ if and only if
$g+(-k)>h$, i.e.\ $g>h+k$. The remaining case $h+k=g$ would give
$(2h+k)^2=4h^2+4hk+k^2=4hg+k^2=D$, again forcing $D$ to be a square. Hence exactly one
of $(h,k)$, $\iota(h,k)$ lies in $P^{+}(D)$, and pairing the elements of $P(D)$ by
$\iota$ gives $\#P^{+}(D)=\frac12\#P(D)$.
\end{proof}

\begin{proof}[Second proof of Theorem~\ref{thm:main}]
Take $D=N^2-4$, which is positive, $\equiv0,1\pmod4$, and not a square for $N\ge3$. By
\eqref{eq:P} and the substitution $m=N-2a$ of Remark~\ref{rem:technau2},
\[
  \#P(D)=\sum_{\substack{k^2<D\\ k^2\equiv D (4)}}d\Bigl(\frac{D-k^2}{4}\Bigr)
        =\sum_{a=1}^{N-1}d\bigl(a(N-a)-1\bigr)=\Phi(N)
\]
by Proposition~\ref{prop:fibre}. By Lemma~\ref{lem:param} and Lemma~\ref{lem:invol},
$Z(D)=\#P^{+}(D)=\frac12\#P(D)=\frac12\Phi(N)$.
\end{proof}

\section{Necklaces: the structural proof}\label{sec:necklace}

We now explain the factor $2$.

\subsection{Necklaces and conjugacy classes}
By Proposition~\ref{prop:free} a matrix $X\in\Mon$ is the same thing as a finite word
$w(X)$ in the letters $\Ll,\Rr$. Since $\tr(XY)=\tr(YX)$, the trace is insensitive to
the starting point of the word: writing $w=uv$ and cutting elsewhere gives
$vu=u^{-1}(uv)u$, a conjugate matrix. The correct objects are therefore \emph{necklaces},
i.e.\ words taken up to cyclic rotation.

\begin{proposition}\label{prop:necklace}
For $N\ge3$, the map $X\mapsto w(X)$ induces a bijection between the
$\SL_2(\Z)$-conjugacy classes of trace $N$ and the necklaces whose matrix has trace $N$.
If $\mathcal{C}$ is such a class and $u_\mathcal{C}$ is the primitive root of the
corresponding necklace, then $\mathcal{C}\cap\Mon$ has exactly
$|u_\mathcal{C}|=\#\Ll(u_\mathcal{C})+\#\Rr(u_\mathcal{C})$ elements. Consequently
\begin{equation}\label{eq:phiLR}
  \Phi(N)=\sum_{\mathcal{C}}\bigl(\#\Ll(u_\mathcal{C})+\#\Rr(u_\mathcal{C})\bigr),
\end{equation}
the sum being over the conjugacy classes of trace $N$.
\end{proposition}

\begin{proof}
Every $X\in\SL_2(\Z)$ with $\tr X\ge3$ is hyperbolic and is conjugate to an element of
$\Mon$, uniquely up to cyclic rotation of its word: this is the classical theory of
cutting sequences and of cycles of reduced indefinite forms, see Series \cite{Series} or
Zagier \cite[\S13]{Zagier}. Rotations of a word $w$ of primitive period $p$ produce
exactly $p$ distinct words, hence, by freeness, exactly $p$ distinct matrices.
\end{proof}

\subsection{The Zagier cycle counts one letter}
Zagier's reduction theory \cite[\S13]{Zagier} attaches to each $\SL_2(\Z)$-class of
indefinite forms of non-square discriminant $D$ a single cycle of Zagier-reduced forms;
consecutive members satisfy $A'=C$ and $B'\equiv-B\pmod{2C}$, and the length of the
cycle is the period of the \emph{minus}, or Hirzebruch--Jung, continued fraction of the
associated quadratic irrational. Forms are not assumed primitive here, in accordance
with the definition of $Z(D)$; correspondingly, the classical bijection between
conjugacy classes of trace $N$ and classes of forms of discriminant $N^2-4$ is onto all
such classes. Combining this with the conversion between regular and minus continued
fractions gives the following.

\begin{proposition}\label{prop:zagierR}
Let $\mathcal{C}$ be a conjugacy class of trace $N\ge3$, with primitive necklace word
$u_\mathcal{C}=\Ll^{a_1}\Rr^{a_2}\cdots\Ll^{a_{2k-1}}\Rr^{a_{2k}}$. Then the Zagier cycle
of the corresponding class of forms of discriminant $N^2-4$ has length
$a_2+a_4+\cdots+a_{2k}=\#\Rr(u_\mathcal{C})$. Consequently
\begin{equation}\label{eq:ZsumR}
  Z(N^2-4)=\sum_{\mathcal{C}}\#\Rr(u_\mathcal{C}).
\end{equation}
\end{proposition}

\begin{proof}
The regular continued fraction with purely periodic part $[a_1,a_2,a_3,\dots,a_{2k}]$
converts into the minus continued fraction with period
\[
  \bigl[a_1+1,\ \underbrace{2,\dots,2}_{a_2-1},\ a_3+2,\
  \underbrace{2,\dots,2}_{a_4-1},\ \dots,\ a_{2k-1}+2,\
  \underbrace{2,\dots,2}_{a_{2k}-1}\bigr].
\]
Each odd-indexed partial quotient contributes one entry and each even-indexed $a_{2i}$
contributes $a_{2i}-1$ entries, so the minus period has length
$k+\sum_{i}(a_{2i}-1)=\sum_i a_{2i}$, which is the number of $\Rr$'s in
$u_\mathcal{C}$. Summing over classes and invoking Zagier's cycle theorem gives
\eqref{eq:ZsumR}.
\end{proof}

\subsection{The swap}
Let
\[
  J=\begin{pmatrix}0&1\\1&0\end{pmatrix}\in\GL_2(\Z),\qquad \det J=-1,\qquad J^2=I .
\]

\begin{lemma}\label{lem:swap}
$J\Ll J=\Rr$ and $J\Rr J=\Ll$. The map $\tau(X)=JXJ$ is an involution of $\Mon$ which
preserves the trace and the determinant; on words it exchanges the letters $\Ll$ and
$\Rr$, hence exchanges $\#\Ll$ and $\#\Rr$; and it maps $\SL_2(\Z)$-conjugacy classes to
$\SL_2(\Z)$-conjugacy classes.
\end{lemma}

\begin{proof}
The two matrix identities are immediate, and they show $\tau(\Mon)=\Mon$ with
$\tau(w)=\bar w$, the word with its letters exchanged. Conjugation by any invertible
matrix preserves the trace and determinant. Finally, if $Y=X^{-1}MX$ then
$\tau(Y)=(JXJ)^{-1}\tau(M)(JXJ)$ and $JXJ\in\SL_2(\Z)$ whenever $X\in\SL_2(\Z)$, since
$\det(JXJ)=\det X$. Thus $\tau$ carries conjugacy classes to conjugacy classes.
\end{proof}

\begin{remark}
It is worth emphasising why this is not a triviality. $J$ has determinant $-1$, so
$\tau$ is an \emph{outer} operation: it is not conjugation by an element of
$\SL_2(\Z)$, and it does not fix conjugacy classes in general --- indeed for $N=4$ it
exchanges the only two classes. What survives is exactly what is needed: it preserves
the trace, so it permutes the finite set of classes of trace $N$, and it exchanges the
two letter counts.
\end{remark}

\begin{proof}[Third proof of Theorem~\ref{thm:main}]
Fix $N\ge3$ and let $\mathcal{S}$ be the (finite) set of conjugacy classes of trace $N$.
By Lemma~\ref{lem:swap}, $\tau$ induces a permutation of $\mathcal{S}$ with
$\#\Ll(u_{\tau(\mathcal{C})})=\#\Rr(u_{\mathcal{C}})$. Summing over $\mathcal{S}$ and
reindexing by $\tau$,
\[
  \sum_{\mathcal{C}\in\mathcal{S}}\#\Ll(u_\mathcal{C})
  =\sum_{\mathcal{C}\in\mathcal{S}}\#\Ll(u_{\tau(\mathcal{C})})
  =\sum_{\mathcal{C}\in\mathcal{S}}\#\Rr(u_\mathcal{C}).
\]
Therefore, by \eqref{eq:phiLR} and \eqref{eq:ZsumR},
\[
  \Phi(N)=\sum_{\mathcal C}\bigl(\#\Ll+\#\Rr\bigr)
        =2\sum_{\mathcal C}\#\Rr
        =2Z(N^2-4). \qedhere
\]
\end{proof}

\begin{remark}
Proposition~\ref{prop:zagierR} is strictly stronger than Theorem~\ref{thm:main}: it is a
statement about each class separately, of which the theorem is the total. The smallest
instance is $N=4$, where the two classes are $\Ll\Ll\Rr$ and $\Ll\Rr\Rr$, exchanged by
$\tau$; their $\Rr$-counts are $1$ and $2$, and discriminant $12$ has exactly
$1+2=3$ Zagier-reduced forms, while $\Phi(4)=6$.
\end{remark}

\section{Necklaces fixed by the swap}\label{sec:swapinv}

For a word $u$ write $\bar u$ for the word with $\Ll$ and $\Rr$ exchanged.

\begin{proposition}\label{prop:swapinv}
Let $w$ be a necklace of length $n$ whose matrix $M$ has trace $N\ge3$. The following
are equivalent:
\begin{enumerate}
\item[(i)] $w$ is fixed by the swap, i.e.\ $\bar w$ is a rotation of $w$;
\item[(ii)] $n$ is even and $w=u\bar u$ as a necklace, for some word $u$ of length $n/2$;
\item[(iii)] $M=V^2$ for some $V\in\GL_2(\Z)$ with $\det V=-1$.
\end{enumerate}
When these hold, $N=s^2+2$ with $s=\tr V\in\Z$. Conversely, for every $s\ge1$ the
necklace $\Ll^s\Rr^s$ is swap-invariant of trace $s^2+2$; hence the set of traces of
swap-invariant necklaces is exactly $\{s^2+2:s\ge1\}$.
\end{proposition}

\begin{proof}
(i)$\Rightarrow$(ii). Let $\sigma$ denote rotation by one letter and suppose
$\bar w=\sigma^j(w)$. Applying the bar operation again and using that it commutes with
$\sigma$ gives $w=\sigma^{2j}(w)$, so $n\mid 2j$ and $j\in\{0,n/2\}$. The case $j=0$
would give $w=\bar w$ letterwise, forcing each letter to equal its own image under the
swap, which is absurd. Hence $n$ is even and $\bar w=\sigma^{n/2}(w)$: writing $w=uv$
with $|u|=|v|=n/2$, this reads $\bar u\bar v=vu$, whence $v=\bar u$.

(ii)$\Rightarrow$(iii). If $w=u\bar u$ then, with $U$ the matrix of $u$, the matrix of
$\bar u$ is $JUJ$ by Lemma~\ref{lem:swap}, so
$M=U\cdot JUJ=(UJ)(UJ)=V^2$ with $V=UJ$ and $\det V=\det U\det J=-1$.

(iii)$\Rightarrow$(i). If $M=V^2$ with $\det V=-1$ then $JMJ=(JVJ)^2$; moreover
$V$ commutes with $M=V^2$, so $M$ and $JMJ$ are conjugate by $VJ\in\SL_2(\Z)$ --- indeed
$(VJ)^{-1}M(VJ)=JV^{-1}V^2VJ=JV^2J=JMJ$ --- and by Proposition~\ref{prop:necklace}
conjugate elements of $\Mon$ have rotation-equivalent words, i.e.\ $\bar w$ is a
rotation of $w$.

For the trace: $\det V=-1$ gives $V^2=(\tr V)V+I$ by Cayley--Hamilton, so
$N=\tr(V^2)=(\tr V)^2+2$. Finally $\Ll^s\Rr^s$ has matrix $\sqm{1+s^2}{s}{s}{1}$ of
trace $s^2+2$, and equals $u\bar u$ with $u=\Ll^s$.
\end{proof}

\section{The cut-sum of a necklace}\label{sec:cutsum}

The necklace point of view makes available an invariant that a linear word does not
possess.

\begin{definition}\label{def:cutsum}
Let $w$ be a necklace of length $n$, and let $M_0,\dots,M_{n-1}\in\Mon$ be the matrices
of its $n$ cuts. The \emph{cut-sum} of $w$ is
\[
  S(w)\;=\;\sum_{i=0}^{n-1}M_i\ \in\ M_2(\Z).
\]
\end{definition}

Rotating the necklace permutes the cuts, so $S(w)$ depends only on the necklace: it is a
genuine matrix-valued invariant, and hence an invariant of the conjugacy class together
with the choice of word length $n$. It refines the trace, which sees only
$\tr S(w)=nN$.

\begin{proposition}\label{prop:detS}
Let $w$ be a necklace of length $n$ and trace $N$. Then $\tr S(w)=nN$ and
\begin{equation}\label{eq:detS}
  \det S(w)\;=\;n+\sum_{0\le i<j\le n-1}\tr\bigl[u_{ij},v_{ij}\bigr],
\end{equation}
where $u_{ij}$ and $v_{ij}$ are the matrices of the two arcs into which the cuts $i$ and
$j$ divide the necklace, and $[u,v]=uvu^{-1}v^{-1}$.
\end{proposition}

\begin{proof}
The trace statement is clear since all $M_i$ have trace $N$. On $M_2$ the determinant is
a quadratic form whose polarisation is
$\det(X+Y)-\det X-\det Y=\tr X\tr Y-\tr(XY)$; applying this repeatedly,
\[
  \det\Bigl(\sum_i M_i\Bigr)=\sum_i\det M_i+\sum_{i<j}\bigl(\tr M_i\tr M_j-\tr(M_iM_j)\bigr)
  = n+\sum_{i<j}\bigl(N^2-\tr(M_iM_j)\bigr).
\]
Fix $i<j$ and let $u=u_{ij}$, $v=v_{ij}$, so that $M_i=uv$ and $M_j=vu$ and
$M_iM_j=uv\,vu$, whence $\tr(M_iM_j)=\tr(u^2v^2)$ by cyclicity. For $X\in\SL_2$,
Cayley--Hamilton gives $X^2=(\tr X)X-I$; writing $x=\tr u$, $y=\tr v$, $z=\tr(uv)=N$,
\[
  \tr(u^2v^2)=\tr\bigl((xu-I)(yv-I)\bigr)=xy\,z-x^2-y^2+2 .
\]
Hence $N^2-\tr(M_iM_j)=x^2+y^2+z^2-xyz-2$, which is $\tr[u,v]$ by Fricke's identity for
$\SL_2$.
\end{proof}

\begin{proposition}\label{prop:minkowski}
Let $w$ be a necklace of length $n$ and trace $N$, put $D=N^2-4$, and set
\[
  \Delta(w)=(\tr S)^2-4\det S,\qquad V(w)=2S(w)-nN\,I .
\]
Then $V(w)$ is an integral traceless matrix and $\Delta(w)=q\bigl(V(w)\bigr)$, where
$q(X)=-\det X$ is the quadratic form of signature $(2,1)$ on the traceless plane. Writing
$M_i=\frac{N}{2}I+P_i$, the traceless parts $P_i$ satisfy $q(P_i)=D/4$ and
\[
  \Delta(w)=4\,q\Bigl(\sum_{i=0}^{n-1}P_i\Bigr).
\]
\end{proposition}

\begin{proof}
$\tr V=2nN-2nN=0$ and $V$ is integral. For traceless $Y$ one has
$\det(\alpha I+Y)=\alpha^2+\det Y$, so with $\alpha=nN/2$ and $Y=\sum P_i$ we get
$\det S=(nN/2)^2+\det\sum P_i$, whence
$\Delta=(nN)^2-4\det S=-4\det\sum P_i=4q(\sum P_i)$. Also $V=2\sum P_i$, so
$q(V)=4q(\sum P_i)=\Delta$. Finally
$q(P_i)=-\det(M_i-\tfrac N2 I)=-\bigl(1-\tfrac{N}{2}\cdot N+\tfrac{N^2}{4}\bigr)
=\tfrac{N^2}{4}-1=\tfrac D4$.
\end{proof}

Geometrically, $P_i$ is the pole vector of the axis of $M_i$, so $\Delta(w)$ is four
times the squared Minkowski length of the sum of the pole vectors of the $n$ lifts of
the closed geodesic determined by $\mathcal{C}$ --- one lift for each crossing of the
Farey tessellation. It is thus a centre-of-mass invariant of a closed geodesic on the
modular surface.

By construction $\Delta$ is invariant under rotation, under reversal of the word, and
under the swap $\tau$; the following was found by exhaustive computation and we state it
as an observation rather than a theorem.

\begin{observation}\label{obs:sep}
Let two necklaces be called \emph{equivalent} when one is obtained from the other by
rotation, reversal, and/or the swap.
\begin{enumerate}
\item[(a)] For every $N\le34$, $\Delta$ takes distinct values on the inequivalent
necklaces of trace $N$; likewise for every $n\le14$, $\Delta$ takes distinct values on
the inequivalent necklaces of length $n$ having equal trace.
\item[(b)] The first coincidences occur at $n=15$, with two pairs, and at $n=16$, with
three pairs. In each of these five cases the entire matrix $S$ agrees, not merely
$\Delta$. For instance, at $n=15$ and $N=576$ the inequivalent necklaces
$\Ll\Ll\Ll\Rr\Ll\Ll\Rr\Rr\Ll\Ll\Ll\Rr\Rr\Ll\Rr$ and
$\Ll\Ll\Ll\Rr\Ll\Ll\Rr\Rr\Ll\Rr\Ll\Ll\Ll\Rr\Rr$ both have
$S=\sqm{4320}{2984}{5176}{4320}$.
\end{enumerate}
\end{observation}

We have no explanation for these coincidences; see \S\ref{sec:open}.

\section{A total-matrix identity}\label{sec:total}

Proposition~\ref{prop:fibre} evaluates not only the cardinality of a trace fibre but any
sum of matrix entries over it.

\begin{theorem}\label{thm:total}
For every $N\ge3$, with $J=\sqm{0}{1}{1}{0}$,
\[
  \sum_{\substack{X\in\Mon\\ \tr X=N}}X
  \;=\;\frac{N}{2}\,\Phi(N)\cdot I\;+\;
  \Bigl(\sum_{a=1}^{N-1}\sigma\bigl(a(N-a)-1\bigr)\Bigr)\cdot J .
\]
In particular $N\Phi(N)$ is even, which also follows from Theorem~\ref{thm:main}.
\end{theorem}

\begin{proof}
By Proposition~\ref{prop:fibre} the fibre is parametrised by pairs $(a,b)$ with
$1\le a\le N-1$ and $b\mid \Nn_a=a(N-a)-1$, the matrix being
$\sqm{a}{b}{\Nn_a/b}{N-a}$. Summing the $(1,1)$ entries gives
$\sum_a a\,d(\Nn_a)$; since $\Nn_a=\Nn_{N-a}$, replacing $a$ by $N-a$ shows this equals
$\sum_a (N-a)d(\Nn_a)$, so twice it is $N\sum_a d(\Nn_a)=N\Phi(N)$ by \eqref{eq:Phidiv}.
The $(2,2)$ entries give the same value. Summing the $(1,2)$ entries gives
$\sum_a\sum_{b\mid \Nn_a}b=\sum_a\sigma(\Nn_a)$, and the $(2,1)$ entries give
$\sum_a\sum_{b\mid \Nn_a}\Nn_a/b$, which is the same because $b\mapsto \Nn_a/b$ permutes
the divisors of $\Nn_a$.
\end{proof}

\begin{remark}
The diagonal of Theorem~\ref{thm:total} is $\Phi$ itself, the object of
\cite{KO,KOPS,Boca,Ustinov,Technau}, whose analysis rests on the divisor sums of Hooley
\cite{Hooley} and Bykovski\u\i--Ustinov \cite{BU}. The off-diagonal entry
$\sum_{a}\sigma(a(N-a)-1)$ is the corresponding sum-of-divisors problem for the same
quadratic polynomial. We are not aware of any treatment of it, and its asymptotic
behaviour --- in particular, whether it is as irregular as $\Phi$ --- appears to be open.
\end{remark}

\section{Computational verification}\label{sec:verify}

All computations were carried out in exact integer arithmetic; no floating-point value
enters any decision (inequalities such as $2h+k>\sqrt D$ are tested as
$2h+k>0$ together with $(2h+k)^2>D$). Both sides of each identity were computed by
independent routines: the left-hand sides by enumerating the free monoid or the
necklaces directly, the right-hand sides by enumerating triples $(A,B,C)$ subject to
\eqref{eq:zagred} and, separately, by the divisor parametrisation of
Lemma~\ref{lem:param}; the resulting values of $\Phi$ and $Z$ were checked against the
terms published in \texttt{A264598} and \texttt{A257007}.

\begin{center}\small
\begin{tabular}{@{}llc@{}}
\toprule
Assertion & Range verified & Outcome\\
\midrule
$\Phi(N)=\sum_a d(a(N-a)-1)=\#P(N^2-4)=2Z(N^2-4)$ & $N=3,\dots,150$ & $148/148$\\
$\sum_\mathcal{C}\#\Ll=\sum_\mathcal{C}\#\Rr=Z(N^2-4)$ & $N=3,\dots,150$ & $148/148$\\
$\iota$ fixed-point-free, condition-reversing, every pair & $N=3,\dots,90$ & no exception\\
Per class: Zagier cycle lengths $=$ $\#\Rr$ as multisets & $N=3,\dots,24$ & identical\\
$\det S=n+\sum\tr[u,v]$ and $\Delta=q(2S-nNI)$ & all necklaces, $n\le9$ & $0$ failures\\
$\Delta$ separates inequivalent necklaces & $N\le34$; $n\le14$ & no collision\\
First $\Delta$-collisions & $n=15,16$ & $2$, then $3$ pairs\\
Theorem~\ref{thm:total} & $N=3,\dots,90$ & $88/88$\\
\bottomrule
\end{tabular}
\end{center}

\medskip
The first ten values are as follows.

\begin{center}
\begin{tabular}{@{}lrrrrrrrrrr@{}}
\toprule
$N$ & 3 & 4 & 5 & 6 & 7 & 8 & 9 & 10 & 11 & 12\\
\midrule
$\Phi(N)$ & 2 & 6 & 8 & 14 & 14 & 24 & 16 & 40 & 26 & 36\\
$Z(N^2-4)$ & 1 & 3 & 4 & 7 & 7 & 12 & 8 & 20 & 13 & 18\\
$\sum_a\sigma(a(N-a)-1)$ & 2 & 10 & 20 & 45 & 62 & 120 & 120 & 270 & 250 & 380\\
\bottomrule
\end{tabular}
\end{center}

\section{Open problems}\label{sec:open}

\begin{enumerate}
\item Explain Observation~\ref{obs:sep}(b). The five known coincidences are presumably
the beginning of an infinite family; is there a construction producing pairs of
inequivalent necklaces with equal cut-sum? Equivalently: which pairs of closed geodesics
of the same length have the same centre of mass over their Farey crossings?
\item Determine the asymptotic behaviour of $\sum_{a=1}^{N-1}\sigma(a(N-a)-1)$, the
off-diagonal entry of Theorem~\ref{thm:total}. Does it admit a limiting distribution
after normalisation, as $\Phi$ does by Peter's theorem \cite{Peter}?
\item Is $\Delta$, or the full cut-sum $S$, expressible in terms of classical invariants
of the corresponding ideal class? Proposition~\ref{prop:minkowski} shows $\Delta$ is the
value at an explicit lattice vector of the ternary form $q$; the configuration of the
$h(N^2-4)$ vectors $V(w)$ attached to a single discriminant deserves study.
\item The OEIS entries \texttt{A264597}, \texttt{A264598} and \texttt{A257007} should
have their conjecture annotations replaced by a reference.
\end{enumerate}

\section*{Data and code availability}

The verification programs are plain Python~3 using only the standard library. Running
\texttt{final.py} performs the end-to-end check of every identity above and reports
either \texttt{ALL CLAIMS VERIFIED} or the first counterexample.

\begin{thebibliography}{99}

\bibitem{Boca} F.~P. Boca,
\emph{Products of matrices $\sqm{1}{1}{0}{1}$ and $\sqm{1}{0}{1}{1}$ and the
distribution of reduced quadratic irrationals},
J. Reine Angew. Math. \textbf{606} (2007), 149--165.

\bibitem{BU} V.~A. Bykovski\u\i\ and A.~V. Ustinov,
\emph{Hooley's problem on representations of a number as the sum of a square and a
product}, Dokl. Math. \textbf{99} (2019), no.~2, 195--200.

\bibitem{Hooley} C.~Hooley,
\emph{On the representation of a number as the sum of a square and a product},
Math. Z. \textbf{69} (1958), 211--227.

\bibitem{KOPS} J.~Kallies, A.~\"Ozl\"uk, M.~Peter and C.~Snyder,
\emph{On asymptotic properties of a number theoretic function arising out of a spin
chain model in statistical mechanics},
Comm. Math. Phys. \textbf{222} (2001), no.~1, 9--43.

\bibitem{KO} P.~Kleban and A.~E. \"Ozl\"uk,
\emph{A Farey fraction spin chain},
Comm. Math. Phys. \textbf{203} (1999), no.~3, 635--647.

\bibitem{Nathanson} M.~B. Nathanson,
\emph{Pairs of matrices in $\GL_2(\mathbb{R}_{\ge0})$ that freely generate},
Amer. Math. Monthly \textbf{122} (2015), no.~8, 790--792.

\bibitem{OEISA264598} OEIS Foundation Inc.,
\emph{Entry \texttt{A264598}}, The On-Line Encyclopedia of Integer Sequences,
\url{https://oeis.org/A264598}. See also \texttt{A264597}.

\bibitem{OEISA257007} OEIS Foundation Inc.,
\emph{Entry \texttt{A257007}: Number of Zagier-reduced binary quadratic forms of
discriminant $n^2-4$}, \url{https://oeis.org/A257007}.

\bibitem{Peter} M.~Peter,
\emph{The limit distribution of a number theoretic function arising from a problem in
statistical mechanics}, J. Number Theory \textbf{90} (2001), no.~2, 265--280.

\bibitem{Series} C.~Series,
\emph{The modular surface and continued fractions},
J. London Math. Soc. (2) \textbf{31} (1985), no.~1, 69--80.

\bibitem{Technau} M.~Technau,
\emph{Remark on the Farey fraction spin chain},
arXiv:2304.08143 [math.NT], 2023.

\bibitem{Ustinov} A.~V. Ustinov,
\emph{Spin chains and Arnold's problem on the Gauss--Kuz'min statistics for quadratic
irrationals}, Sb. Math. \textbf{204} (2013), no.~5, 762--779.

\bibitem{Zagier} D.~B. Zagier,
\emph{Zetafunktionen und quadratische K\"orper}, Springer, Berlin, 1981.

\end{thebibliography}

\end{document}
